%% ---------------------------------------------------------------------
\section{Análisis, sistema y control}
\label{sec:in-palabras}
%% ---------------------------------------------------------------------

El nombre del curso tiene tres palabras y cada una carga con un significado preciso. Vale la pena detenerse en ellas antes de escribir la primera ecuación.

\paragraph{Análisis}
El diccionario de la Real Academia da varias acepciones de \emph{análisis}: la distinción y separación de las partes de algo para conocer su composición, el estudio detallado de algo y, la que nos interesa, la \emph{parte de las matemáticas basada en los conceptos de límite, convergencia y continuidad}, que da origen al cálculo diferencial e integral y a la teoría de funciones. En el curso las dos ideas van juntas: \emph{separamos} un sistema en sus partes (planta, actuador, sensor, controlador) y \emph{estudiamos} cada una con herramientas del análisis matemático: ecuaciones diferenciales, transformada de Laplace, variable compleja.

\paragraph{Sistema}
Un sistema es un conjunto de cosas (personas, células, moléculas, tuberías, transistores) interconectadas de tal forma que producen un patrón de comportamiento a medida que pasa el tiempo; puede ser activado, restringido o impulsado por fuerzas externas, pero su respuesta depende de sus propias características~\citep{MeWr08}. Más corto: una combinación de componentes que se percibe como un todo~\citep{HaCh08}. Lo que hace interesantes a los sistemas que estudiaremos es que son \emph{dinámicos}: lo que hacen ahora depende de lo que les pasó antes. En el capítulo~\ref{cap:po} precisaremos estas ideas con las nociones de señal y de sistema dinámico.

\paragraph{Control}
La Real Academia también da varias acepciones de \emph{control} (del francés \emph{contrôle}): comprobación, dominio, el examen parcial con el que se comprueba la marcha de los alumnos\ldots{} y la que nos interesa: \emph{regulación, manual o automática, sobre un sistema}.

\begin{definicion}[label=def:in-sistemacontrol]{Sistema de control}
Un \clave{sistema de control} es un conjunto de dispositivos que administra, ordena, dirige o regula el comportamiento de otro sistema para obtener un resultado deseado~\citep{WikiControl}. Está formado por subsistemas (el proceso o planta, los actuadores, los sensores, el controlador) ensamblados para que una salida siga una referencia o rechace perturbaciones, con un desempeño especificado.
\end{definicion}

El \emph{análisis de sistemas de control}, entonces, estudia cómo evolucionan en el tiempo esos elementos interconectados y qué tan bien cumplen su propósito: si el lazo es estable, qué tan rápido y con cuánto sobrepaso responde, cuánto error deja, qué tanto le afectan las perturbaciones y los errores del modelo.

\subsection{La idea central: retroalimentar}

\begin{definicion}[label=def:in-retro]{Retroalimentación}
\clave{Retroalimentar} es medir la variable que se quiere controlar, compararla con el valor deseado y usar la diferencia (el \emph{error}) para decidir la acción. Un sistema así se llama \clave{retroalimentado} o \emph{de lazo cerrado}.
\end{definicion}

La retroalimentación es el concepto fundamental del curso y su gran virtud es la \clave{robustez}: como la acción depende de lo que realmente pasa, el sistema puede cumplir sus requisitos aunque el modelo sea imperfecto, aparezcan perturbaciones o las mediciones lleguen con ruido. Su peligro es la \alerta{sobrecorrección}: si se reacciona demasiado fuerte o demasiado tarde, el lazo oscila y puede volverse inestable. Casi todo el curso trata de cómo aprovechar la virtud sin caer en el peligro.

Richard Murray resumió el área en su conferencia Bode de 2016 con una ecuación de cuatro términos, \emph{control = dynamics, uncertainty, feedforward, feedback}~\citep{Mur16}, cuyas iniciales (DUFF) recuerdan la cerveza de \emph{Los Simpson}. Leída en orden: (i) hay unas \clave{dinámicas} que queremos modificar; (ii) las conocemos con \clave{incertidumbre}; (iii) con la \clave{prealimentación} (\emph{feedforward}) actuamos con lo que sabemos de antemano: la referencia, o una perturbación que podemos medir; y (iv) con la \clave{retroalimentación} (\emph{feedback}) corregimos lo que no sabíamos.

\begin{ideaclave}
\textbf{Señales y sistemas.}
La retroalimentación se entiende con dos conceptos entrelazados~\citep{AlMa10}: las \emph{señales}, porque es midiéndolas como se decide qué hacer, y los \emph{sistemas}, porque la retroalimentación existe para cambiar su comportamiento. Cuando hablamos de un sistema hablamos de algo que produce o modifica una señal; cuando hablamos de una señal, de algo que algún sistema produjo.
\end{ideaclave}

\begin{pausa}
Cuando usted conduce un carro por una carretera con curvas, ¿quién hace de sensor, de controlador y de actuador? ¿Qué es la referencia y qué la perturbación?
\respuesta{Los ojos (y el oído interno) son el sensor; el cerebro compara la posición del carro con el centro del carril (la referencia) y decide: es el controlador; los brazos y el timón son el actuador. El viento lateral, los baches y la inclinación de la vía son perturbaciones. Si usted cierra los ojos, el lazo se abre.}
\end{pausa}
